\documentclass[aspectratio=169,11pt]{beamer}

\usetheme{Madrid}
\usecolortheme{default}
\setbeamertemplate{navigation symbols}{}
\setbeamertemplate{footline}[frame number]

\usepackage[T1]{fontenc}
\usepackage[utf8]{inputenc}
\usepackage{amsmath,amssymb}
\usepackage{listings}
\usepackage{xcolor}
\usepackage{hyperref}

\definecolor{codegray}{RGB}{245,245,245}
\definecolor{codeblue}{RGB}{36,82,145}
\definecolor{codegreen}{RGB}{40,120,60}

\lstset{
  basicstyle=\ttfamily\scriptsize,
  keywordstyle=\color{codeblue},
  commentstyle=\color{gray},
  stringstyle=\color{codegreen},
  breaklines=true,
  frame=single,
  backgroundcolor=\color{codegray},
  columns=fullflexible,
  keepspaces=true
}

\title{Parsing Python with Haskell and Lowering to MLIR}
\author{Jegannarayanan Sivakasiulaganathan}
\date{\today}

\begin{document}

% ------------------------------------------------------------
\begin{frame}
	\titlepage
	\vspace{-0.4cm}
	\begin{center}
		\small \url{https://github.com/handel-fan/efp_final}
	\end{center}
\end{frame}

% ------------------------------------------------------------
\begin{frame}{Project Goal}
	\begin{block}{Goal}
		Build a small compiler-style pipeline that parses Python using Haskell and lowers a restricted Python subset toward MLIR.
	\end{block}

	\vspace{0.3cm}
	\centering
	\Large
	Python source $\rightarrow$ Haskell AST $\rightarrow$ Project AST $\rightarrow$ MLIR AST $\rightarrow$ MLIR

	\vspace{0.5cm}
	\normalsize
	The completed work focuses mainly on parsing and AST conversion. The MLIR lowering step was started but not fully completed.
\end{frame}

% ------------------------------------------------------------
\begin{frame}{What is MLIR?}
	\begin{itemize}
		\item MLIR is an intermediate representation used between high-level source code and
		      lower-level LLVM IR.
		\item It supports multiple levels of abstraction before reaching machine code.
		\item Dialects group related operations and types for specific domains.
	\end{itemize}

	\vspace{0.35cm}
	\centering
	\Large
	Source code $\rightarrow$ MLIR $\rightarrow$ LLVM IR $\rightarrow$ machine code
\end{frame}

% ------------------------------------------------------------
\begin{frame}[fragile]{MLIR Example}
	\begin{lstlisting}
module {
  func.func @add_example(%a: i32, %b: i32) -> i32 {
    %0 = arith.addi %a, %b : i32
    return %0 : i32
  }
}
  \end{lstlisting}

	\vspace{0.25cm}
	\begin{itemize}
		\item \texttt{func} describes functions.
		\item \texttt{arith} describes arithmetic operations.
		\item This form enables machine-agnostic analysis and optimization.
	\end{itemize}
\end{frame}

% ------------------------------------------------------------
\begin{frame}{MLIR Core Concepts}
	\begin{description}
		\item[Operations] Basic units of computation. Operations take inputs and may produce
		      results.
		\item[Regions] Containers attached to operations. They hold blocks and represent
		      bodies of functions or control flow.
		\item[Blocks] Ordered sequences of operations. Blocks may also take arguments.
		\item[Dialects] Namespaces for related operations and types, such as \texttt{arith},
		      \texttt{func}, and \texttt{llvm}.
	\end{description}
\end{frame}

% ------------------------------------------------------------
\begin{frame}{What is an AST?}
	\begin{block}{Abstract Syntax Tree}
		An abstract syntax tree, or AST, is a tree-shaped representation of a program's structure after parsing.
	\end{block}

	\vspace{0.25cm}
	\begin{itemize}
		\item It removes surface-level syntax details such as whitespace and punctuation.
		\item It keeps the important structure of the program: variables, literals,
		      expressions, statements, and functions.
		\item Compiler stages usually operate on this structured representation instead of
		      raw source text.
	\end{itemize}

	\vspace{0.3cm}
	\centering
	\Large
	Source text $\rightarrow$ Parser $\rightarrow$ AST $\rightarrow$ Compiler passes
\end{frame}

% ------------------------------------------------------------
\begin{frame}[fragile]{Example: Python Source to AST}
	\begin{columns}[T,onlytextwidth]
		\begin{column}{0.42\textwidth}
			\textbf{Python input}
			\begin{lstlisting}[language=Python]
x: int = 5
y: int = 10
z: int = x + y
      \end{lstlisting}
		\end{column}

		\begin{column}{0.55\textwidth}
			\textbf{Simplified AST shape}
			\begin{lstlisting}
Module
  Assign x
    Type: int
    Value: IntLit 5

  Assign y
    Type: int
    Value: IntLit 10

  Assign z
    Type: int
    Value:
      BinaryOp +
        Var x
        Var y
      \end{lstlisting}
		\end{column}
	\end{columns}

	\vspace{0.2cm}
	The AST represents the meaning-bearing structure of the program, not the exact text formatting.
\end{frame}

% ------------------------------------------------------------
\begin{frame}{The Converter: Main Project Work}
	\begin{block}{Main idea}
		The bulk of the project was converting the large \texttt{language-python} AST into a smaller project-defined AST.
	\end{block}
	\centering
	\Large
	\texttt{language-python AST}
	$\rightarrow$
	\texttt{ConvertAST.hs}
	$\rightarrow$
	\texttt{PMMAST.hs AST}
\end{frame}
\begin{frame}{Why Use Haskell to Parse Python?}
	\begin{itemize}
		\item Haskell represents syntax directly with algebraic data types.
		\item Expressions, statements, and functions can be modeled explicitly instead of
		      using generic objects or loosely structured values.
		\item Parser and converter logic can be composed from smaller functions.
		\item Types such as \texttt{Either} and \texttt{Maybe} make failure cases visible in
		      function signatures.
		\item Compared with C++, Haskell avoids much of the boilerplate around classes,
		      pointers, references, ownership, and error handling.
	\end{itemize}
\end{frame}

% ------------------------------------------------------------
\begin{frame}{PythonMLIR-- Language Subset}
	\begin{itemize}
		\item The project uses a restricted Python subset to keep the compiler pipeline
		      manageable.
		\item The subset focuses on typed variables, integer operations, expressions, and
		      basic control flow.
		\item Dynamic Python features are intentionally avoided.
	\end{itemize}

	\vspace{0.35cm}
	\begin{block}{Planning documents}
		\texttt{docs/planning/PythonMLIR--.md}: supported and excluded features.\newline
		\texttt{docs/planning/pythonmlir\_supported\_ops.py}: concrete supported syntax examples.
	\end{block}
\end{frame}

% ------------------------------------------------------------
\begin{frame}[fragile]{Example Lowering Target}
	\begin{columns}[T,onlytextwidth]
		\begin{column}{0.47\textwidth}
			\textbf{PythonMLIR-- input}
			\begin{lstlisting}[language=Python]
x: int = 5
y: int = 10
z: int = x + y
      \end{lstlisting}
		\end{column}
		\begin{column}{0.50\textwidth}
			\textbf{Simplified MLIR output}
			\begin{lstlisting}
%0 = arith.constant 5 : i32
%1 = arith.constant 10 : i32
%2 = arith.addi %0, %1 : i32
      \end{lstlisting}
		\end{column}
	\end{columns}

	\vspace{0.35cm}
	The intended lowering maps typed integer declarations and expressions into MLIR arithmetic operations.
\end{frame}

% ------------------------------------------------------------
\begin{frame}{What Works}
	\begin{itemize}
		\item A defined \texttt{PythonMLIR--} subset.
		\item Concrete examples of supported syntax and operations.
		\item Conversion from the \texttt{language-python} AST to a smaller project AST.
		\item A project AST designed around the selected Python features.
		\item A partial lowering path from the project AST toward MLIR structures.
	\end{itemize}
\end{frame}

% ------------------------------------------------------------
\begin{frame}{What I Learned}
	\begin{itemize}
		\item How to convert and walk ASTs using Haskell.
		\item How host AST structure determines what the converter must handle.
		\item Haskell constructs such as monads and record syntax.
		\item MLIR concepts including operations, regions, blocks, terminators, block
		      arguments, and control flow.
		\item Compiler projects become difficult quickly because each stage depends on
		      precise structure from the previous stage.
	\end{itemize}
\end{frame}

% ------------------------------------------------------------
\begin{frame}{Conclusion}
	\begin{block}{Main result}
		The project successfully scoped a Python subset and implemented the front half of the pipeline: parsing Python and converting it into a smaller Haskell-defined AST.
	\end{block}

	\vspace{0.3cm}
	\begin{block}{Next step}
		Complete \texttt{LowerMLIR.hs} so the project AST lowers into \texttt{mlir-hs} structures and emits MLIR code automatically.
	\end{block}
\end{frame}

\end{document}
